\begin{textsample}{NEG Dim 2 – 2013\_c5 – Score -2.00 – 2013\_c5/enade\_2013\_fisio\_a31223fe5a.txt}  \label{ex:f2_neg_063}
CAPÍTULO 5 DISTRIBUIÇÃO DOS CONCEITOS 5.1 PANORAMA NACIONAL DA DISTRIBUIÇÃO DOS CONCEITOS A Tabela 5.1 apresenta a quantidade e distribuição de cursos de Fisioterapia participantes do ENADE/2013, por faixa de conceito e Grande Região.
A diferença entre os cursos tabulados neste capítulo e no capítulo 2 são os cursos sem conceito, em princípio, aqueles sem alunos concluintes que participassem da prova ou que tivessem tão somente um aluno concluinte¹³.
Observando -se os dados da Tabela 5.1, nota -se que, dos 372 cursos participantes, 148 ( 39,8 \% ) classificaram -se com conceito 3, o valor modal.
Este foi também o conceito modal em todas as regiões do Brasil ( entre 34,1 \% na região Sudeste e 50,6 \% na região Nordeste ), exceto na região Centro-Oeste, onde a moda foi o conceito 4 ( 32,4 \% ).
O conceito 4 foi o segundo mais frequente em nível nacional ( 26,3 \%, correspondendo a 98 cursos ) e o conceito 2, o terceiro ( 22,0 \%, correspondendo a 82 cursos ).
Houve, ainda, 15 cursos ( 4,0 \% ) que receberam conceito 5 e outros dez cursos que receberam conceito 1.
Dos 372 cursos de Fisioterapia, 19 ficaram sem conceito ( SC ).
A região Norte participou com 22 cursos ou 5,9 \% do total nacional.
Nesta região, o conceito 3 foi atribuído a dez cursos dos 22 participantes, o que equivale a 45,5 \% do total regional.
Apenas um curso ( 4,5 \% ) ficou sem conceito ( SC ).
Os demais foram avaliados com conceito 2 ( nove cursos correspondendo a 40,9 \% ) e conceito 4 ( dois cursos correspondendo a 9,1 \% ).
Nenhum curso foi alocado aos conceitos 1 ou 5. ¹³ Estes cursos com somente um concluinte não foram considerados no capítulo 2.
A região Nordeste participou com 79 cursos ou 21,2 \% do total nacional.
Nesta região, 40 cursos, 50,6 \% em termos regionais, obtiveram conceito 3, o conceito modal para a região.
Os conceitos 2 e 4 foram atribuídos a, respectivamente, 20 e 12 cursos ( 25,3 \% e 15,2 \% ).
Já os conceitos 1 e 5 foram atribuídos a, respectivamente, três e quatro cursos ( 3,8 \% e 5,1 \% ).
Nessa região, nenhum curso ficou sem conceito ( SC ).
Como já comentado, dos 170 cursos participantes da região Sudeste, 58 ( 34,1 \% ) obtiveram conceito 3, o conceito modal.
O conceito 1 foi atribuído a cinco cursos ( 2,9 \% ) e o conceito 2, a 37 ( 21,8 \% ).
O conceito 4 foi atribuído a 48 cursos ( 28,2 \% ) e seis outros ( 3,5 \% ) receberam o conceito 5.
Enquanto 16 cursos ( 9,4 \% ) ficaram sem conceito ( SC ).
A região Sul também contou com cursos \textbf{distribuídos} em todas as faixas de conceitos.
A predominância do conceito 3 foi de 45,3 \%, como já comentado anteriormente, correspondentes a 29 dos 64 cursos participantes na região Sul.
O conceito 2 foi atribuído a seis cursos ( 9,4 \% ) e o conceito 4, a 24 cursos ( 37,5 \% ).
O conceito 1 foi recebido por um único curso ( 1,6 \% ) e o conceito 5, por quatro deles ( 6,3 \% ).
Nenhum dos cursos da região Sul ficou sem conceito ( SC ).
Como já comentado, de todos os 37 cursos participantes na região Centro-Oeste, 12 receberam conceito 4, o conceito modal ( 32,4 \% ).
Nesta região apenas dois cursos ficaram sem conceito ( SC ).
Já, um curso ( 2,7 \% ) foi alocado no conceito 1 e outro único ( 2,7 \% ) curso no conceito 5, correspondendo ao total de 5,4 \% dos cursos da região.
Os demais cursos foram avaliados com conceito 2 ( dez cursos, 27,0 \% ) e conceito 3 ( 11 cursos, 32,4 \% ).
O Gráfico 5.1 apresenta a distribuição do Conceito ENADE segundo Grande Região : poligonais mais à esquerda representam Grandes Regiões com pior distribuição, e poligonais mais à direita, Grandes Regiões com melhores conceitos.
A região Norte ( linha azul ) apresentou os cursos com valores mais baixos, ao passo que a região Sul ( linha roxa ) apresentou a distribuição com valores maiores.
Já os cursos das regiões Nordeste ( linha verde ), Sudeste ( linha vermelha ) e Centro-Oeste ( linha preta ) apresentaram situação intermediária. 5.2 CONCEITOS POR CATEGORIA ADMINISTRATIVA E POR GRANDE REGIÃO A Tabela 5.2 apresenta a distribuição dos cursos participantes do ENADE/2013 de Fisioterapia, por Categoria Administrativa, de acordo com os conceitos por eles alcançados, segundo as Grandes Regiões brasileiras.
Dos 372 cursos participantes, 48 ( 12,9 \% ) eram ministrados em Instituições Públicas e 324 ( 87,1 \% ), em Privadas.
De acordo com as informações da Tabela 5.2, em termos nacionais, das 15 IES que receberam o conceito 5, 13 eram Públicas e apenas duas Privadas.
Dos 48 cursos participantes de IES Públicas, o conceito 4 foi o valor modal, atribuído a 26 cursos ( 54,2 \% ), enquanto nenhum curso recebeu conceito 1.
Entre os demais cursos de Instituições Públicas participantes, dois obtiveram conceito 2 ( 4,2 \% da categoria ), cinco receberam conceito 3 ( 10,4 \% ) e como já comentado, os 13 restantes foram alocados no conceito 5 ( 27,1 \% ).
Nesta categoria, apenas dois cursos ficaram sem conceito ( SC ).
Na rede privada, o conceito modal foi 3, com 143 cursos dos 324 da categoria.
Entre os demais cursos participantes, dez receberam conceito 1 e 80, conceito 2.
O conceito 4 foi atribuído a 72 cursos e o conceito 5, como já comentado, somente a dois cursos.
Nesta categoria administrativa, 17 dos cursos ficaram sem conceito ( SC ).
Na análise por região, observa -se que, na região Norte, as Instituições Privadas participaram com 19 cursos ( 86,4 \% do total regional ), dos quais apenas um ficou sem conceito ( SC ).
O conceito modal para as Instituições Privadas na região foi 3, com nove cursos, correspondendo a 47,4 \%.
O conceito 2 apresenta frequência muito próxima do conceito modal, com oito cursos, correspondendo a 42,1 \%.
Apenas um curso foi alocado no conceito 4 ( 5,3 \% ), enquanto nenhum curso foi alocado nos conceitos 1 ou 5.
As Instituições Públicas participaram com três cursos ( 13,6 \% do total regional ), dos quais nenhum ficou sem conceito ( SC ) ou recebeu os conceitos 1 ou 5 nessa categoria.
Os três cursos ficaram alocados igualmente aos conceitos 2, 3 e 4 ( um curso, cada ).
Na região Nordeste, a rede privada concentrou 67 dos 79 cursos participantes, equivalentes a 84,8 \% do total da região.
Nesta Categoria Administrativa na região, \textbf{predominaram} os cursos com conceito 3, 40 cursos correspondendo a 59,7 \% dos cursos oferecidos por IES Privadas no Nordeste.
Apenas três cursos receberam o conceito 1 ( 4,5 \% ), enquanto o conceito 2 participou com 20 cursos ( 29,9 \% ) e o conceito 4 foi avaliado com quatro cursos correspondendo a 6,0 \%.
Nenhum curso ficou sem conceito ( SC ) ou recebeu conceito 5 nessa categoria e região.
As Instituições Públicas dessa região participaram com 12 cursos ( 15,2 \% ), dos quais oito obtiveram conceito 4, o valor modal, seguidos de quatro cursos que receberam conceito 5.
Nenhum curso ficou sem conceito ( SC ), ou recebeu os conceitos 1, 2 ou 3, nesta combinação de categoria e região.
Na região Sudeste, a proporção de cursos da rede privada, 89,4 \%, foi mais elevada do que nas demais regiões brasileiras, correspondendo a 152 dos 170 cursos participantes.
Nesta categoria, na região Sudeste, o conceito modal foi 3 ( 56 cursos ).
Nesta combinação de Categoria Administrativa e Grande Região, 14 dos cursos não receberam conceito ( SC ).
Os demais foram avaliados com conceito 1 ( cinco cursos ), conceito 2 ( 36 cursos ), conceito 4 ( 39 cursos ) e conceito 5 ( dois cursos ).
Entre os 18 cursos em Instituições Públicas na região Sudeste, a categoria modal foi a que recebeu conceito 4 ( nove cursos ).
Os demais foram avaliados com os conceitos 2 ( um curso ), 3 ( dois cursos ) e 5 ( quatro cursos ).
Na região Sudeste, nesta categoria, apenas dois cursos ficaram sem conceito ( SC ), enquanto nenhum curso recebeu conceito 1.
As Instituições Privadas concentraram 53 dos 64 cursos participantes da região Sul, 82,8 \% do total regional, a menor proporção entre todas as regiões.
Desses, 28 ficaram com conceito 3, o conceito modal.
Nesta combinação de Categoria Administrativa e Grande Região, nenhum curso ficou sem receber conceito ( SC ) ou recebeu conceito 5.
Os demais foram avaliados com conceito 1 ( um curso ), conceito 2 ( seis cursos ) e conceito 4 ( 18 cursos ).
As Instituições Públicas na região Sul participaram com 11 cursos ( 17,2 \% ), dos quais nenhum ficou sem conceito ( SC ) ou recebeu os conceitos 1 ou 2.
O conceito modal foi 4, com seis cursos.
Além disso, os demais foram avaliados com conceito 3 ( um curso ) e conceito 5 ( quatro cursos ).
Na região Centro-Oeste, 33 dos 37 cursos participantes eram de Instituições Privadas ( 89,2 \% em termos regionais ).
Destes, 30 cursos ficaram alocados igualmente aos conceitos modais 2, 3 e 4 ( dez cursos, cada ).
Além destes, somente um curso recebeu o conceito 1, enquanto apenas dois cursos ficaram sem conceito ( SC ) nesta categoria.
Dos quatro cursos de Instituições Públicas, dois foram avaliados no conceito 4, conceito modal.
Outros dois cursos receberam igualmente os conceitos 3 e 5 ( um curso, cada ).
Nesta região, nenhum curso de IES Pública ficou sem conceito ( SC ).
O Gráfico 5.2 apresenta a distribuição do Conceito ENADE segundo a Categoria Administrativa da IES.
Os cursos em IES Públicas ( linha azul ) apresentam uma poligonal mais à direita do que os cursos das IES Privadas ( linha verde ) e consequentemente uma distribuição de Conceitos ENADE com valores maiores.
Nota -se, também, que para IES com conceito 1 não houve diferença entre o tipo de Categoria Administrativa. 5.3 CONCEITOS POR ORGANIZAÇÃO ACADÊMICA E POR GRANDE REGIÃO Na Tabela 5.3 encontra -se a distribuição dos conceitos atribuídos aos cursos participantes do ENADE/2013 na Área de Fisioterapia, por Organização Acadêmica, segundo as Grandes Regiões brasileiras.
Dos 372 cursos participantes, 158 eram oferecidos em Universidades, 152 em Faculdades e os demais 62 em Centros Universitários.
Esta distribuição corresponde a, respectivamente, 42,5 \%, 40,9 \% e 16,7 \% dos cursos.
De acordo com os dados apresentados, todos os 15 cursos avaliados com conceito 5 eram vinculados a Universidades.
Apenas cinco cursos em Universidades ficaram sem conceito ( SC ).
Esse tipo de Organização Acadêmica teve o conceito 4 como modal, com 60 cursos ( 38,0 \% ).
Os demais cursos de Universidades avaliados receberam os conceitos 1 ( três cursos ), 2 ( 26 cursos ), 3 ( 49 cursos ) e conceito 5 ( 15 cursos, como já mencionado ).
Entre os cursos em Centros Universitários, o conceito modal foi 3, com 26 cursos ( 41,9 \% ).
Neste tipo de organização acadêmica nenhum curso ficou com conceito 5.
Os outros cursos neste tipo de Organização Acadêmica receberam os conceitos 1 ( dois cursos ), 2 ( 11 cursos ) e 4 ( 19 cursos ).
Apenas quatro cursos ficaram sem conceito ( SC ).
Nas Faculdades, dez dos 152 cursos ficaram sem conceito ( SC ) e 73 ( 48,0 \% ) receberam o mesmo conceito modal dos Centros Universitários, conceito 3.
Dos demais cursos neste tipo de Organização Acadêmica, cinco receberam conceito 1, 45 conceito 2 e 19 conceito 4.
Nenhum curso recebeu conceito 5 nessa categoria.
Considerando -se separadamente as regiões brasileiras, verifica -se que, na região Norte, as Universidades concentraram 5 dos 22 cursos participantes.
Dos cursos em Universidades, nenhum curso ficou sem conceito ( SC ) ou recebeu conceitos 1 ou 5.
Quatro cursos ficaram alocados igualmente aos conceitos 2 e 3 ( dois cursos, cada ).
Somente um curso recebeu conceito 4.
Os Centros Universitários da região Norte obtiveram o mesmo número de cursos e a mesma distribuição dos conceitos das Universidades na mesma região.
As Faculdades participaram com 12 cursos na região Norte.
Somente um curso ficou sem conceito ( SC ).
O conceito modal foi 3, com seis cursos seguido de cinco cursos alocados ao conceito 2.
Nenhum curso recebeu os demais conceitos restantes nessa categoria e região.
Na região Nordeste, as Universidades participaram com 19 dos 79 cursos na Área de Fisioterapia da região.
Nenhum dos cursos em Universidades no Nordeste ficou sem conceito ( SC ) ou recebeu os conceitos 1 ou 2.
O conceito modal foi 4, com oito cursos.
Os demais receberam conceito 3 ( sete cursos ) e conceito 5 ( quatro cursos ).
Os Centros Universitários contaram com oito cursos participantes na região Nordeste, sendo quatro deles com conceito 2, o conceito modal, e os quatro restantes com conceitos 3 e 4 ( dois cursos, cada ).
Nenhum curso ficou sem conceito ( SC ) ou recebeu conceitos 1 ou 5.
As Faculdades foram representadas por 52 cursos na região Nordeste, dos quais apenas três ficaram com conceito 1.
O conceito modal foi 3, com 31 cursos.
Dos cursos restantes, 16 receberam conceito 2 e dois conceito 4.
Nenhum curso ficou sem conceito ( SC ) ou recebeu conceito 5.
Na região Sudeste, as Universidades concentraram 87 dos 170 cursos da região.
Entre os cursos em Universidades na região, o conceito modal foi 4 com 29 cursos, e apenas cinco cursos ficaram sem conceito ( SC ).
Os demais cursos receberam os conceitos 1(dois cursos ), 2 ( 20 cursos ), 3 ( 25 cursos ) e 5 ( seis cursos ).
Os Centros Universitários participaram com 35 cursos na região Sudeste, dos quais 16 obtiveram conceito modal, 3, e apenas quatro ficaram sem conceito ( SC ).
Os demais receberam os conceitos 1 ( dois cursos ), 2 ( três cursos ) e 4 ( dez cursos ).
Nenhum curso recebeu conceito 5 nessa categoria e região.
As Faculdades foram representadas por 48 cursos na região Sudeste, que se \textbf{distribuíram} nos conceitos 1 ( um curso ), 2 ( 14 cursos ), 3 ( 17 cursos, conceito modal ) e 4 ( nove cursos ).
Apenas sete cursos ficaram sem conceito ( SC ), enquanto nenhum recebeu o conceito 5 nessa combinação de categoria e região.
Dos 64 cursos da região Sul, 36 eram de Universidades, para os quais o conceito modal foi 4, com 16 cursos.
Nesse tipo de organização, nenhum dos cursos ficou sem conceito ( SC ).
Os demais receberam os conceitos 1 ( um curso ), 2 ( três cursos ), 3 ( 12 cursos ) e 5 ( quatro cursos ).
Os Centros Universitários da região Sul tiveram cinco dos oito cursos participantes no conceito modal, 3.
Os três cursos restantes receberam conceito 4, enquanto nenhum curso ficou sem conceito ( SC ) ou recebeu os demais conceitos restantes nessa categoria e região.
Foram 20 os cursos vinculados a Faculdades na região Sul, e nenhum destes ficou Sem conceito ( SC ) ou recebeu os conceitos 1 ou 5.
De todos os cursos, 12 receberam conceito 3, o conceito modal.
Os demais receberam os conceitos 2 ( três cursos ) e 4 ( cinco cursos ).
Na região Centro-Oeste, 11 dos 37 cursos eram de Universidades.
Nesse tipo de organização, nenhum curso ficou sem conceito ( SC ) ou recebeu conceito 1.
O conceito modal foi 4, com seis cursos.
Os outros cursos obtiveram os conceitos 2 ( um curso ), 3 ( três cursos ) e 5 ( também, um curso ).
Os Centros Universitários da região Centro-Oeste contaram com seis cursos, dos quais dois receberam o conceito 2, um o conceito 3 e quatro cursos o conceito 4, o conceito modal.
Nenhum curso ficou sem conceito ( SC ) ou recebeu os conceitos 1 ou 5.
Dos 20 cursos em Faculdades na região Centro-Oeste, nenhum ficou com conceito 5, um recebeu conceito 1, sete conceitos 2, outros sete o conceito 3 e três conceito 4.
Apenas dois cursos ficaram sem conceito ( SC ) nessa categoria e região.
O Gráfico 5.3 apresenta a distribuição do Conceito ENADE segundo a Organização Acadêmica da IES.
Os cursos em Universidades ( linha azul ) apresentam uma poligonal mais à direita do que as demais, denotando uma distribuição de Conceitos ENADE com valores maiores.
No outro extremo temos os cursos em Faculdades ( linha vermelha ) com a poligonal mais à esquerda e consequentemente uma distribuição de Conceitos ENADE com menores valores.
IES de Centros Universitários ( linha verde ) apresentam uma situação intermediária.

% matched lemmas: distribuir, predominar
\end{textsample}
